\chapter{C++ moderne : lambdas, déplacement, possession}

Le chapitre sur le C++ donnait le langage ; celui sur les templates, la
généricité. Il reste ce qui sépare un programme qui compile d'un programme qu'on
ose modifier : les \textbf{lambdas}, la \textbf{sémantique de déplacement}, et
les objets qui \textbf{possèdent} ce qu'ils tiennent.

Ces trois sujets sont ceux qu'on repousse en croyant qu'ils sont avancés. Ils ne
le sont pas : ils sont \emph{quotidiens}, et les ignorer produit du code qui
copie sans le savoir et qui fuit sans le dire.

\section{Les lambdas : une fonction écrite là où on s'en sert}

\begin{nkcode}{Syntaxe}
auto doubler = [](int32 x) { return x * 2; };
int32 y = doubler(21);                        // 42
\end{nkcode}

Une lambda est un objet-fonction que le compilateur fabrique pour vous. Les
crochets sont la \textbf{capture} : ce que la lambda emporte de son entourage.

\begin{center}
\begin{tabular}{@{}ll@{}}
\toprule
\textbf{Capture} & \textbf{Sens} \\
\midrule
\texttt{[]} & rien --- la lambda ne voit que ses paramètres \\
\texttt{[x]} & une \emph{copie} de \texttt{x}, figée à la création \\
\texttt{[\&x]} & une \emph{référence} à \texttt{x} --- il doit vivre plus longtemps \\
\texttt{[=]} & copie de tout ce qui est employé \\
\texttt{[\&]} & référence sur tout ce qui est employé \\
\texttt{[this]} & l'objet courant, par référence \\
\bottomrule
\end{tabular}
\end{center}

\begin{nkretenir}
\textbf{Le dépôt s'en sert dès la première ligne d'un programme.} Chaque
application déclare ses données ainsi :

\begin{nkcode}{Applications/NkDames/src/main.cpp}
NKENTSEU_DEFINE_APP_DATA(([]() {
    NkAppData d;
    d.appName = "NkDames";
    // ...
    return d;
}()));
\end{nkcode}

Une lambda \emph{appelée sur place} --- notez le \texttt{()} final. Elle sert à
construire un objet en plusieurs étapes et à le rendre \textbf{constant} :
plusieurs affectations, un seul résultat, et pas de variable temporaire qui
traîne dans la portée englobante.
\end{nkretenir}

\begin{nkpiege}[Le piège qui coûte le plus cher]
\textbf{Une capture par référence sur un objet qui meurt avant la lambda est un
pointeur pendant, avec un autre nom.}

\begin{nkcode}{FAUX}
NkFunctionCallback FabriquerRappel() {
    int32 compteur = 0;
    return [&compteur]() { return ++compteur; };   // compteur MEURT ici
}
\end{nkcode}

Le programme compile, et fonctionne souvent --- la pile n'a pas encore été
réécrite. C'est exactement le défaut du chapitre sur le C, sous un déguisement
qui ne ressemble plus du tout à un pointeur.

\textbf{La règle} : une lambda qui \emph{survit} à la portée où elle est écrite
--- un rappel, un gestionnaire d'événement, une tâche donnée à un pool --- capture
\textbf{par valeur}. \texttt{[\&]} ne va qu'avec un usage immédiat.
\end{nkpiege}

\begin{nkretenir}
\texttt{[this]} capture le \emph{pointeur} sur l'objet, jamais une copie. Une
lambda stockée dans un membre et capturant \texttt{this} devient invalide quand
l'objet meurt --- ce qui arrive plus souvent qu'on ne croit, dans une interface
où les panneaux naissent et disparaissent.
\end{nkretenir}

\section{Le déplacement : transférer au lieu de copier}

Un vecteur d'un million d'éléments passé par valeur se copie : un million de
constructions et une allocation. Or le plus souvent la source ne sert plus après.

\begin{nkcode}{Kernel/Foundation/NKCore/src/NKCore/NkTraits.h}
template <typename T> constexpr NkRemoveReference_t<T> &&NkMove(T &&value) noexcept {
    return static_cast<NkRemoveReference_t<T> &&>(value);
}
\end{nkcode}

\begin{nkretenir}
\textbf{\texttt{NkMove} ne déplace rien.} C'est une conversion de type : elle dit
« traite cette valeur comme temporaire ». Le déplacement effectif est fait par le
\emph{constructeur de déplacement} de la classe, qui a le droit de voler les
entrailles de la source puisqu'elle est réputée finie.

C'est pourquoi compter les \texttt{NkMove} d'un fichier ne mesure aucune
allocation : cela mesure une \emph{intention}.
\end{nkretenir}

\begin{nkcode}{Ce que doit faire un constructeur de deplacement}
NkTampon(NkTampon &&autre) noexcept
    : mData(autre.mData), mTaille(autre.mTaille) {
    autre.mData = nullptr;      // INDISPENSABLE
    autre.mTaille = 0;
}
\end{nkcode}

\begin{nkpiege}
\textbf{Oublier de vider la source est la faute classique, et elle est
silencieuse.} Sans les deux dernières lignes, \emph{deux} objets tiennent le
même pointeur ; le premier destructeur libère, le second libère à nouveau. Une
double libération corrompt le tas --- et le plantage arrive plus tard, ailleurs,
dans du code innocent.

\textbf{La source doit rester détruisible, pas utilisable.} C'est le contrat
exact : après un déplacement, on n'a le droit que de détruire l'objet ou de lui
affecter une nouvelle valeur.
\end{nkpiege}

\begin{nkretenir}
\textbf{\texttt{noexcept} n'est pas décoratif sur un déplacement.} Un conteneur
qui s'agrandit ne déplacera ses éléments que si leur déplacement est déclaré
\texttt{noexcept} ; sinon il les \emph{copie}, pour pouvoir revenir en arrière en
cas d'exception.

Un mot-clé oublié, et la réallocation copie au lieu de déplacer --- sans qu'aucun
message ne le dise, et sans que le résultat soit faux. Seulement lent.
\end{nkretenir}

\section{La possession : qui libère, et quand}

Le C++ moderne répond à « qui libère ? » par un \emph{type}. Nkentseu en offre
quatre, tous dans \texttt{NKMemory} et tous en en-tête seul.

\begin{center}
\begin{tabular}{@{}lll@{}}
\toprule
\textbf{Type} & \textbf{Possession} & \textbf{Quand} \\
\midrule
\texttt{NkUniquePtr} & un seul propriétaire & le cas par défaut \\
\texttt{NkSharedPtr} & plusieurs, compteur atomique & durée de vie vraiment partagée \\
\texttt{NkWeakPtr} & aucune --- observe & briser un cycle \\
\texttt{NkIntrusivePtr} & compteur \emph{dans} l'objet & entités, composants, ressources \\
\bottomrule
\end{tabular}
\end{center}

\begin{nkretenir}
\textbf{Commencez toujours par \texttt{NkUniquePtr}.} Il ne coûte rien de plus
qu'un pointeur nu, il dit qui libère, et il libère tout seul. On ne passe à
\texttt{NkSharedPtr} que si la durée de vie est \emph{réellement} partagée ---
c'est-à-dire si l'on ne peut pas nommer un propriétaire.

« Je ne sais pas qui doit libérer » n'est pas une raison de partager : c'est le
signe qu'une décision de conception n'a pas été prise.
\end{nkretenir}

\begin{nkretenir}
\texttt{NkIntrusivePtr} range le compteur \emph{dans} l'objet, ce qui économise
une allocation par rapport à \texttt{NkSharedPtr} et son bloc de contrôle
séparé. En échange, l'objet doit être conçu pour --- on ne peut pas l'appliquer à
un type qu'on ne contrôle pas.

L'en-tête le destine explicitement aux « objets qui DOIVENT être référencés :
Entity, Component, Resource ». C'est un choix guidé par le nombre d'objets : à
cent mille entités, une allocation de moins par entité se compte.
\end{nkretenir}

\begin{nkpiege}[Le cycle]
Deux objets qui se tiennent l'un l'autre par \texttt{NkSharedPtr} ne meurent
\textbf{jamais} : chacun maintient le compteur de l'autre à un. Ce n'est pas une
fuite banale --- aucun outil ne signale une mémoire dont le compteur est
légitimement positif.

La parade est \texttt{NkWeakPtr} : le lien « retour » observe sans posséder. La
règle de conception qui l'accompagne : \textbf{dans une relation
parent/enfant, le parent possède, l'enfant observe.}
\end{nkpiege}

\section{Deux types qui remplacent des conventions}

\begin{nkcode}{Kernel/Foundation/NKCore/src/NKCore/NkOptional.h et NkVariant.h}
NkOptional<int32> Chercher(...);   // « peut-etre un entier »
NkVariant<A, B, C> valeur;         // « exactement un des trois »
\end{nkcode}

\begin{nkretenir}
\textbf{\texttt{NkOptional} remplace la valeur sentinelle}, celle qu'on choisit
faute de mieux : \texttt{-1}, \texttt{0}, \texttt{nullptr}, la chaîne vide. Une
sentinelle a deux défauts que le type supprime : elle est une \emph{valeur
possible} du domaine (le jour où $-1$ devient légitime, tout casse), et rien
n'oblige l'appelant à la tester.

\texttt{NkVariant} remplace l'union plus un champ de type --- même contenu, mais
c'est le compilateur qui garantit qu'on lit le bon membre, au lieu d'une
convention que chacun doit respecter.
\end{nkretenir}

\section{Le reste du confort, en une page}

\begin{center}
\begin{tabular}{@{}ll@{}}
\toprule
\textbf{Écriture} & \textbf{Ce qu'elle apporte} \\
\midrule
\texttt{auto x = f();} & le type suit sa source quand elle change \\
\texttt{for (auto \&e : v)} & parcours sans indice --- donc sans erreur d'indice \\
\texttt{constexpr} & calculé à la compilation, coût nul à l'exécution \\
\texttt{nullptr} & un pointeur nul \emph{typé}, contrairement à \texttt{0} \\
\texttt{enum class} & pas de conversion implicite en entier \\
\texttt{= delete} & interdit explicitement une opération \\
\texttt{override} & le compilateur vérifie que vous redéfinissez vraiment \\
\bottomrule
\end{tabular}
\end{center}

\begin{nkpiege}
\textbf{\texttt{for (auto e : v)} copie chaque élément.} Sur un vecteur de
chaînes, c'est une allocation par tour, et rien ne le dit.

Écrivez \texttt{auto \&} pour modifier, \texttt{const auto \&} pour lire. Le
\texttt{auto} nu ne se justifie que sur des types minuscules --- un entier, un
flottant.
\end{nkpiege}

\begin{nkretenir}
\textbf{\texttt{override} est le mot-clé au meilleur rendement du langage.} Sans
lui, une signature légèrement différente de celle de la base --- un \texttt{const}
oublié, un type de paramètre voisin --- crée une \emph{nouvelle} fonction au lieu
d'en redéfinir une. Tout compile, et la méthode de la base continue d'être
appelée.

Avec \texttt{override}, le compilateur refuse. Une erreur au lieu d'un
comportement.
\end{nkretenir}

\begin{nkbilan}
Une lambda emporte ce qu'elle capture : par valeur si elle survit à sa portée,
par référence seulement pour un usage immédiat. \texttt{NkMove} est une
conversion, pas un déplacement ; le vrai travail est dans le constructeur de
déplacement, qui doit \textbf{vider la source} et se déclarer \texttt{noexcept}.
La possession se dit par un type : \texttt{NkUniquePtr} par défaut,
\texttt{NkSharedPtr} quand la durée de vie est vraiment partagée,
\texttt{NkWeakPtr} pour briser les cycles, \texttt{NkIntrusivePtr} quand les
objets se comptent par milliers. Enfin \texttt{NkOptional} et \texttt{NkVariant}
remplacent deux conventions que rien ne faisait respecter.
\end{nkbilan}

\begin{nkquestions}
\begin{enumerate}
\item Quand faut-il capturer par valeur, et quand par référence ?
\item \texttt{NkMove} déplace-t-il quelque chose ? Que fait-il exactement ?
\item Que doit faire un constructeur de déplacement à la source, et que
      se passe-t-il s'il l'oublie ?
\item Pourquoi \texttt{noexcept} change-t-il le comportement d'un conteneur qui
      s'agrandit ?
\item Quand préférer \texttt{NkIntrusivePtr} à \texttt{NkSharedPtr} ?
\end{enumerate}
\end{nkquestions}

\begin{nkpratique}
Écrivez une classe \texttt{NkTampon} qui possède un tableau d'octets alloué par
NKMemory. Donnez-lui les cinq opérations : destructeur, constructeur de copie,
affectation par copie, constructeur de déplacement, affectation par déplacement.

Puis remplissez un \texttt{NkVector<NkTampon>} de mille éléments, et affichez le
nombre de copies et de déplacements réellement effectués --- un compteur statique
incrémenté dans chaque constructeur suffit.

Retirez ensuite le \texttt{noexcept} du constructeur de déplacement, et refaites
la mesure.
\end{nkpratique}

\begin{nkdemo}
Prouvez qu'un cycle de \texttt{NkSharedPtr} ne se libère jamais.

Créez deux objets qui se tiennent mutuellement par \texttt{NkSharedPtr}, chacun
affichant un message dans son destructeur. Sortez de la portée : \textbf{aucun
message}. Remplacez l'un des deux liens par \texttt{NkWeakPtr} : les deux
messages apparaissent.

\textbf{Ce que la démonstration doit établir}, et c'est la partie qu'on saute :
ne vous contentez pas de constater l'absence de message. Affichez aussi le
\emph{compteur de références} des deux objets juste avant la sortie de portée.
Sans ce chiffre, « rien ne s'affiche » pourrait aussi vouloir dire que votre
destructeur n'affiche pas.
\end{nkdemo}

\begin{nklogique}
Une fonction range des lambdas dans une liste de rappels, appelée plus tard.
Certains rappels donnent des résultats absurdes, de façon intermittente.

\begin{enumerate}
\item Quelle est la cause la plus probable ? Décrivez le mécanisme précis.
\item Pourquoi le défaut est-il \emph{intermittent} plutôt que systématique ?
      Votre réponse doit expliquer pourquoi il disparaît souvent sous le
      débogueur.
\item Un collègue propose de tout capturer par valeur, partout, comme règle
      générale. Donnez le cas où cette règle est coûteuse, et le cas où elle est
      \emph{fausse} --- car il en existe un.
\end{enumerate}
\end{nklogique}
